% from ega2-1-fr.tex
\oldpage[II]{60}

\begin{proof}
La première assertion résulte aussitôt de
(\hyperref[II.3.4.4-fr]{3.4.4}) en prenant
$\mathcal{M}=\Gamma_*(\mathcal{F})$. Pour établir la seconde, il suffira
d'établir que $\mathcal{M}$ est limite inductive de ses sous-$\mathcal{S}$-Modules
gradués de type fini $\mathcal{N}_\lambda$~: en effet, il en résultera que
$\widetilde{\mathcal{M}}$ est limite inductive des
$\widetilde{\mathcal{N}}_\lambda$
(\hyperref[II.3.2.4-fr]{3.2.4}), donc $\mathcal{F}$ est limite inductive des
$\beta(\widetilde{\mathcal{N}}_\lambda)$~; comme $X$ est quasi-compact
((\hyperref[II.3.4.1-fr]{3.4.1}) et
\hyperref[I.6.3.1-fr]{I, 6.3.1}) et que $\mathcal{F}$ est de type fini,
$\mathcal{F}$ sera nécessairement égal à un des
$\beta(\widetilde{\mathcal{N}}_\lambda)$
(\hyperref[0.5.2.3-fr]{0, 5.2.3}).

Pour définir les $\mathcal{N}_\lambda$ ayant $\mathcal{M}$ pour limite
inductive, il suffit de considérer pour chaque $n\in\mathbf{Z}$ le
$\mathcal{O}_Y$-Module quasi-cohérent $\mathcal{M}_n$, qui est limite
inductive de ses sous-$\mathcal{O}_Y$-Modules
$\mathcal{M}_n^{(\mu_n)}$ de type fini, en raison des hypothèses sur $Y$
(\hyperref[I.9.4.9-fr]{I, 9.4.9})~; il est immédiat que
$\mathcal{P}_{\mu_n}=\mathcal{S}\cdot\mathcal{M}_n^{(\mu_n)}$ est un
$\mathcal{S}$-Module gradué de type fini, et on vérifie aussitôt qu'en prenant
pour $\mathcal{N}_\lambda$ les sommes finies de $\mathcal{S}$-Modules de la
forme $\mathcal{P}_{\mu_n}$, on répond bien à la question.
\end{proof}

\begin{corollary}[3.4.6]
\label{II.3.4.6-fr}
Supposons vérifiées les hypothèses de (\hyperref[II.3.4.4-fr]{3.4.4}), et en
outre que l'espace sous-jacent $Y$ soit quasi-compact~; soit $\mathcal{F}$ un
$\mathcal{O}_X$-Module quasi-cohérent de type fini. Il existe $n_0$ tel que
pour $n\geq n_0$, l'homomorphisme canonique
$\sigma:p^*(p_*(\mathcal{F}(n)))\to\mathcal{F}(n)$
(\hyperref[0.4.4.3-fr]{0, 4.4.3}) soit surjectif.
\end{corollary}

\begin{proof}
En effet, pour tout $y\in Y$, soit $U$ un voisinage ouvert affine de $y$ dans
$Y$. Il existe un entier $n_0(U)$ tel que, pour $n\geq n_0(U)$,
$\mathcal{F}(n)|p^{-1}(U)$ soit engendré par un nombre fini de ses sections
au-dessus de $p^{-1}(U)$ (\hyperref[II.2.7.9-fr]{2.7.9})~; mais ces dernières
sont images canoniques de sections de $p^*(p_*(\mathcal{F}(n)))$ au-dessus de
$p^{-1}(U)$ (\hyperref[0.3.7.1-fr]{0, 3.7.1} et
\hyperref[0.4.4.3-fr]{0, 4.4.3}), donc
$\mathcal{F}(n)|p^{-1}(U)$ est égal à l'image canonique de
$p^*(p_*(\mathcal{F}(n)))|p^{-1}(U)$. Enfin, comme $Y$ est quasi-compact, il
existe un recouvrement fini de $Y$ par des ouverts affines $U_i$, et en
prenant pour $n_0$ le plus grand des $n_0(U_i)$, on achève la démonstration.
\end{proof}

\par\medskip\noindent\textit{Remarques (3.4.7). ---\ }
\label{II.3.4.7-fr}
Si $p=(\psi,\theta):X\to Y$ est un morphisme d'espaces annelés et
$\mathcal{F}$ un $\mathcal{O}_X$-Module, le fait que l'homomorphisme canonique
$\sigma:p^*(p_*(\mathcal{F}))\to\mathcal{F}$ soit surjectif s'explicite de la
façon suivante (\hyperref[0.4.4.1-fr]{0, 4.4.1})~: pour tout $x\in X$ et toute
section $s$ de $\mathcal{F}$ au-dessus d'un voisinage ouvert $V$ de $x$, il
existe un voisinage ouvert $U$ de $p(x)$ dans $Y$, un nombre fini de sections
$t_i$ ($1\leq i\leq m$) de $\mathcal{F}$ au-dessus de $p^{-1}(U)$, un
voisinage $W\subset V\cap p^{-1}(U)$ de $x$ et des sections $a_i$
($1\leq i\leq m$) de $\mathcal{O}_X$ au-dessus de $W$ telles que
\[
  s|W=\sum_i a_i\cdot(t_i|W).
\]
Lorsque $Y$ est un \emph{schéma affine} et $p_*(\mathcal{F})$
\emph{quasi-cohérent}, cette condition équivaut au fait que $\mathcal{F}$ est
\emph{engendré par ses sections au-dessus de $X$}
(\hyperref[0.5.5.1-fr]{0, 5.5.1})~: en effet, si
$Y=\operatorname{Spec}(A)$, on peut supposer que $U=D(f)$ avec $f\in A$~; il
existe un entier $n>0$ et des sections $s_i$ de $\mathcal{F}$ au-dessus de
$X$ telles que $t_i$ soit la restriction à $p^{-1}(U)$ de $s_ig^n$, avec
$g=\theta(f)$ (en appliquant (\hyperref[I.1.4.1-fr]{I, 1.4.1}) à
$p_*(\mathcal{F})$)~; comme $g$ est inversible au-dessus de $p^{-1}(U)$, on a
donc
\[
  s|W=\sum_i b_i\cdot(s_i|W)
\]
avec $b_i=a_i(g|W)^{-n}$, d'où notre assertion. Lorsque $Y$ est affine, le
corollaire (\hyperref[II.3.4.6-fr]{3.4.6}) redonne donc
(\hyperref[II.2.7.9-fr]{2.7.9}).

